\documentclass[12pt,a4paper,oneside]{article}
\usepackage[QX]{polski}

\usepackage[utf8]{inputenc}
\usepackage{latexsym}
\usepackage{tgpagella}
\usepackage{lmodern}
\usepackage{amsmath,amsthm,amsfonts,amssymb,alltt}
\usepackage{epsfig}
\usepackage{pdflscape}
\usepackage{caption}
\usepackage{indentfirst}
\usepackage{float}
%\usepackage{showkeys}
\bibliographystyle{plabbrv}


\usepackage{color}
\usepackage[polish]{babel}
\usepackage{datetime2}
\usepackage[x11names,dvipsnames,table]{xcolor}
\usepackage{hyperref}
\hypersetup{
pdfauthor={Roman Czapla, Olaf Bar},
colorlinks=True,
linkcolor=darkgray,  % color of internal links (change box color with linkbordercolor)
citecolor=BrickRed,  % color of links to bibliography
filecolor=Magenta,   % color of file links
urlcolor=BlueViolet}	%%pdfpagemode=FullScreen}

% diagramy, grafy itp.
\usepackage{tikz}
\usetikzlibrary{positioning}
\usetikzlibrary{arrows}
\usetikzlibrary{arrows.meta}
\usetikzlibrary{chains,fit,shapes,calc}
\tikzset{main node/.style={circle,fill=blue!20,draw,minimum size=1cm,inner sep=0pt}}

% algorytmy
\usepackage[linesnumbered,lined,commentsnumbered]{algorithm2e}
\SetKwFor{ForEach}{for each}{do}{end for}%
\SetKwFor{ForAll}{for all}{do}{end for}%
\newenvironment{myalgorithm}
{\rule{\textwidth}{0.5mm}\\\SetAlCapSty{}\SetAlgoNoEnd\SetAlgoNoLine\begin{algorithm}}{\end{algorithm}\rule{\textwidth}{0.5mm}}


%---------------------
\overfullrule=2mm
\pagestyle{plain}
\textwidth=15cm \textheight=685pt \topmargin=-25pt \linespread{1.3} 
\setlength{\parskip}{0pt}
\setlength\arraycolsep{2pt}
\oddsidemargin =0.9cm
\evensidemargin =-0.1cm

\captionsetup{width=.95\linewidth, justification=centering}
%---------------------




\newtheorem{tw}{Twierdzenie}[section]
\newtheorem{lem}[tw]{Lemat}
\newtheorem{co}[tw]{Wniosek}
\newtheorem{prop}[tw]{Stwierdzenie}
\theoremstyle{definition}
\newtheorem{ex}{Przykład}
\newtheorem{re}[tw]{Uwaga}
\newtheorem{de}{Definicja}[section]



\newcommand{\bC}{{\mathbb C}}
\newcommand{\bR}{{\mathbb R}}
\newcommand{\bZ}{{\mathbb Z}}
\newcommand{\bQ}{{\mathbb Q}}
\newcommand{\bN}{{\mathbb N}}
\newcommand{\captionT}[1]{\caption{\textsc{\footnotesize{#1}}}}
\renewcommand\figurename{Rys.}

\numberwithin{equation}{section}
\renewcommand{\thefootnote}{\arabic{footnote})}
%\renewcommand{\thefootnote}{\alph{footnote})}



\begin{document}

% --------------------------------------------
% Strona tytułowa
% --------------------------------------------

\thispagestyle{empty}
\begin{titlepage}
\begin{center}\Large
Uniwersytet Komisji Edukacji Narodowej w Krakowie\\
\large
Instytut Bezpieczeństwa i Informatyki\\
\vskip 10pt
\end{center}
\begin{center}
\centering \includegraphics[width=1.0\columnwidth]{images/logo.png}
\end{center}

\begin{center}
 {\bf \fontsize{14pt}{14pt}\selectfont PROJEKT INŻYNIERSKI\\ WYBÓR TEMATU}
\end{center}
\vskip 5pt
\begin{center}
 {\bf \fontsize{22pt}{22pt}\selectfont System WasteFree Cloud do efektywnego zarządzania odpadami}
\end{center}

\begin{center}
 {\fontsize{12pt}{12pt}\selectfont wykonany przez: }
\end{center}
\begin{center}
 {\bf\fontsize{16pt}{16pt}\selectfont Dawid Plata}\\
 {\fontsize{12pt}{12pt}\selectfont Nr albumu: 165303 \\\&\\}
 {\bf\fontsize{16pt}{16pt}\selectfont Michał Szymbara}\\
 {\fontsize{12pt}{12pt}\selectfont Nr albumu: 165297 \\\&\\}
 {\bf\fontsize{16pt}{16pt}\selectfont Michał Wechmann}\\
 {\fontsize{12pt}{12pt}\selectfont Nr albumu: 165298}
\end{center}
\begin{center}
 {\fontsize{12pt}{12pt}\selectfont pod opieką:}\\
 {\bf\fontsize{12pt}{12pt}\selectfont dr Wojciech Gwizdała}\\
 {\bf\fontsize{12pt}{12pt}\selectfont mgr Wojciech Baran}
\end{center}

%\mbox{}
\vspace*{\fill}
%\vskip 50pt
\begin{center}
\large
Kraków \the\year
\end{center}
\end{titlepage}
\setcounter{page}{0} 
\newpage\null\thispagestyle{empty}
%\setcounter{page}{0} 
%\newpage
%\thispagestyle{empty}

\tableofcontents


\newpage

\section{Informacje na temat wybranego tematu i zespołu projektowego}
\subsection{Wybrany temat}
Wybranym tematem projektu jest zaprojektowanie i wykonanie systemu \textbf{WasteFree Cloud} --- dwujęzycznego (PL/EN) systemu wspierającego efektywne zarządzanie procesem zamawiania oraz realizacji wywozu odpadów dla użytkowników indywidualnych, grup domowych i operatorów wywozowych.

System obejmuje pełny cykl obsługi: rejestrację i autoryzację, tworzenie grup i adresów, składanie oraz akceptację zleceń, komunikację w czasie rzeczywistym, rozliczenia oraz dokumentację realizacji (w tym rozliczenie opłaty utylizacyjnej).
\subsection{Zespół projektowy}
Zespół projektowy:
\begin{itemize}
\item Dawid Plata (e-mail: dawid.plata@student.up.krakow.pl)
\item Michał Szymbara (e-mail: michal.szymbara@student.up.krakow.pl)
\item Michał Wechmann (e-mail: michal.wechmann@student.up.krakow.pl)
\end{itemize}

Założony podział ról w zespole (wstępny) zakłada pracę równoległą nad warstwą interfejsu (UI), API oraz warstwami domenową i infrastrukturalną, przy wspólnym ustalaniu wymagań i scenariuszy testowych.
\subsection{Opis tematu}
WasteFree Cloud to \textbf{system} wspierający uporządkowanie i usprawnienie procesu wywozu odpadów poprzez ujednolicenie komunikacji pomiędzy trzema stronami: użytkownikami, grupami (np. gospodarstwa domowe / sąsiedztwo) oraz operatorami wywozowymi.

W ramach systemu użytkownik może:
\begin{itemize}
\item założyć konto (z potwierdzeniem e-mail) oraz zalogować się,
\item uzupełnić profil (m.in. opis, avatar, IBAN do wypłat, saldo konta pre-paid),
\item działać w ramach grup: tworzyć grupy, zapraszać członków, przyjmować zaproszenia oraz zarządzać adresem wywozu,
\item tworzyć zlecenia wywozu (typ odpadu, ilość, preferowane przedziały czasowe, notatki) wraz z podziałem kosztów pomiędzy uczestników,
\item akceptować/odrzucać zaproszenia i zlecenia poprzez skrzynkę powiadomień (Inbox) oraz otrzymywać notyfikacje (aplikacja, e-mail, push web),
\item śledzić status zleceń oraz rozliczenia (w tym opłatę utylizacyjną) w ramach historii i powiadomień.
\end{itemize}

Operator wywozowy otrzymuje listę dostępnych zleceń w danym obszarze (np. mieście), może zaakceptować zlecenie do realizacji, dostarczyć dokumentację wykonania oraz wnioskować o wypłatę środków po zatwierdzeniu przez użytkowników.

System jest dwujęzyczny (PL/EN) i wykorzystuje lokalizację komunikatów po stronie API.
\subsection{Założenia projektowe}
Poniżej przedstawiono wstępne założenia funkcjonalne i technologiczne wynikające z koncepcji systemu oraz aktualnej implementacji w repozytorium.

\subsubsection*{Założenia funkcjonalne}
\begin{itemize}
\item Dwujęzyczność interfejsu i komunikatów (język polski i angielski).
\item Obsługa ról i scenariuszy: \emph{Użytkownik/Członek grupy}, \emph{Operator wywozowy} oraz konto administracyjne.
\item Pełny proces biznesowy: rejestracja i aktywacja konta, tworzenie grup i zarządzanie adresem wywozu, składanie zleceń, akceptacje uczestników, przyjęcie zlecenia przez operatora (\emph{garbage admin}), realizacja oraz rozliczenia (w tym dopłaty/zwroty wynikające z opłaty utylizacyjnej).
\item Komunikacja i aktualizacje w czasie rzeczywistym: czat grupowy oraz powiadomienia (SignalR/WebSocket).
\item Rejestrowanie kluczowych zdarzeń i zmian (m.in. pola audytowe dla encji utrwalanych w bazie danych).
\end{itemize}

\subsubsection*{Założenia techniczne (stack)}
\begin{itemize}
\item \textbf{Frontend}: Angular 17 (TypeScript), Bootstrap 5, RxJS; komunikacja z API po HTTP oraz w czasie rzeczywistym przez SignalR (biblioteka \texttt{@microsoft/signalr}). 
\item \textbf{Backend}: ASP.NET (Minimal APIs) na platformie \textbf{.NET}, architektura warstwowa (API / Application / Domain / Infrastructure), obsługa logiki w formie poleceń/zapytań z własnym mechanizmem mediatora, walidacja danych wejściowych przy użyciu FluentValidation.
\item \textbf{Dane}: Entity Framework Core; środowisko lokalne oparte o SQLite.
\item \textbf{Geokodowanie i adresy}: integracja z \textbf{Nominatim API} (OpenStreetMap) do wyszukiwania/adresowania lokalizacji.
\item \textbf{Asynchroniczność i zadania}: mechanizm zadań cyklicznych (TickerQ) dla procesów w tle.
\item \textbf{Storage plików}: składowanie plików użytkowników w zewnętrznym storage (Azure Blob Storage), w tym avatarów oraz dowodu opłaty utylizacyjnej (zdjęcie/skan/PDF) przesyłanego przez operatora (\emph{garbage admin}).
\item \textbf{Uruchamianie i środowiska}: konteneryzacja Docker; lokalne uruchomienie całości przez Docker Compose (oddzielnie usługa API oraz UI).
\end{itemize}

\subsubsection*{Założenia bezpieczeństwa}
\begin{itemize}
\item Uwierzytelnienie oparte o token JWT oraz autoryzację ról/polityk.
\item Potwierdzenie konta poprzez aktywację linkiem wysłanym na e-mail.
\end{itemize}
\subsection{Wstępny harmonogram pracy}

Wstępny harmonogram prac nad aplikacją (ustalenia przed rozpoczęciem realizacji):
\begin{itemize}
\item Analiza wymagań i doprecyzowanie zakresu funkcjonalnego (MVP), ról oraz scenariuszy użycia.
\item Projekt: architektura, model danych, kontrakty API oraz wstępne makiety kluczowych widoków.
\item Przygotowanie środowiska pracy i wdrożeń: Docker/Compose, podstawowe CI/CD oraz konfiguracja środowisk.
\item Implementacja Backend (API, autoryzacja, baza danych, integracje) oraz równolegle Implementacja Frontend (UI, komunikacja z API).
\item Integracja całości (w tym komunikacja czasu rzeczywistego), testy manualne i poprawki.
\item Uzupełnienie dokumentacji oraz finalne przygotowanie wydania.
\end{itemize}

% --------------------------------------------------------------------
%%%%%%% odkomentować gdy bibliografia ma być wewnątrz dokumentu
% --------------------------------------------------------------------
%\begin{thebibliography}{11}
%
%\addcontentsline{toc}{section}{Literatura}
%
%\bibitem{ZAN}
%C. Zannoni and P. Pasini, 
%\emph{Advances in the Computer Simulatons of Liquid Crystals}, Kluwer Academic Publishers, 2000.
%
%\end{thebibliography}

\end{document}

